
\chapter{Concrete Package-Sameness Design Contract}
\label{ch:v154-psame-design-contract}

The v1.5.3 reflection theorem is correct for the one-constructor concrete package-sameness relation. It must not be silently transferred to an equivalence closure. This chapter fixes the active design.

\begin{definition}[Base concrete package-sameness]
\label{def:v154-base-psame}
For a fixed generated probe bundle \(\Pi\), the active concrete signature-token package-sameness relation is written
\[
  \psame_{\Pi}^{\mathrm{base}}(p,q).
\]
It is generated by exactly one constructor:
\[
  \TokIntro(\Pi,s,p),\quad
  \TokIntro(\Pi,t,q),\quad
  \hsame(s,t)
  \quad\Longrightarrow\quad
  \psame_{\Pi}^{\mathrm{base}}(p,q).
\]
The active notation \(\psame_{\Pi}^{\mathrm{sig}}\) is identified with this base relation:
\[
  \psame_{\Pi}^{\mathrm{sig}} := \psame_{\Pi}^{\mathrm{base}}.
\]
\end{definition}

\begin{principle}[No implicit equivalence closure]
\label{prin:v154-no-implicit-closure}
No theorem in the concrete signature Globalize instance may replace \(\psame_{\Pi}^{\mathrm{base}}\) by an equivalence closure without explicitly adding a closure-reflection theorem.
\end{principle}

\begin{definition}[Optional equivalence closure]
\label{def:v154-psame-eq}
If a future package-algebra layer needs an equivalence relation, it must introduce a separate relation
\[
  \psame_{\Pi}^{\mathrm{eq}}(p,q)
  := \EqCl(\psame_{\Pi}^{\mathrm{base}})(p,q).
\]
This relation is not the active relation for v1.5.4 concrete Globalize exactness.
\end{definition}

\begin{definition}[Closure-reflection obligation]
\label{def:v154-closure-reflection}
The optional equivalence-closure relation is reflection-safe only if the following additional theorem is proved:
\[
\begin{aligned}
  &\TokIntro(\Pi,s,p),\quad
    \TokIntro(\Pi,t,q),\quad
    \psame_{\Pi}^{\mathrm{eq}}(p,q)\\
  &\hspace{3.0cm}\Longrightarrow\quad \hsame(s,t).
\end{aligned}
\]
This obligation is named \(\ClosureReflect(\Pi)\).
\end{definition}

\begin{theorem}[Base reflection remains the active theorem]
\label{thm:v154-base-reflection-active}
Assume \(\TokUnique(\Pi)\). If
\[
  \TokIntro(\Pi,s,p),\quad
  \TokIntro(\Pi,t,q),\quad
  \psame_{\Pi}^{\mathrm{base}}(p,q),
\]
then
\[
  \hsame(s,t).
\]
\end{theorem}

\begin{proof}
This is Theorem~\ref{thm:v153-package-reflection-token-unique} with the notation of Definition~\ref{def:v154-base-psame}. The only constructor of \(\psame_{\Pi}^{\mathrm{base}}\) exposes introducing signatures \(s_0,t_0\) with \(\hsame(s_0,t_0)\). Token-introduction uniqueness compares \(s\) with \(s_0\) and \(t\) with \(t_0\); symmetry and transitivity of \(\hsame\) finish the proof.
\end{proof}

\begin{corollary}[Equivalence closure is not exported as exact]
\label{cor:v154-eq-not-exported}
The exact concrete Globalize theorem of v1.5.4 exports \(\psame_{\Pi}^{\mathrm{base}}\), not \(\psame_{\Pi}^{\mathrm{eq}}\). Exporting \(\psame_{\Pi}^{\mathrm{eq}}\) as exact is deferred until \(\ClosureReflect(\Pi)\) is checked.
\end{corollary}

\begin{proof}
By Principle~\ref{prin:v154-no-implicit-closure}, equivalence closure is a separate construction. Definition~\ref{def:v154-closure-reflection} records the extra theorem required for reflection through that closure. Since v1.5.4 records no checked proof of that theorem, only the base relation is exported as exact.
\end{proof}
